\section{Dependent elimination, indices, and transports}\label{sec:dependent}

\subsection{Case analysis is a motive-construction problem}

For ordinary sums, case analysis gives independent-looking branches. For an
indexed family or a dependent local context, it also changes the types of
objects that remain in scope. A generic rule must therefore build a legal
elimination motive, not splice a surface \code{match} around existing holes.

Given a scrutinee $s:I\,\bar p\,\bar i$, the case rule should perform these
operations in one transaction:
\begin{enumerate}
\item Compute the dependency closure of $s$ and its indices in the local
context, goal, and coupled obligations. Generalize the dependent declarations
needed by the chosen elimination operation.
\item Inspect the actual eliminator and its permitted result sort. Construct a
motive over the indices and scrutinee that returns the generalized goal.
\item Instantiate each constructor branch, introducing its exact fields and
index equalities. Reintroduce dependent context entries with their transformed
types and update the branch-to-parent substitution.
\item Generate branches or discharge impossible constructor/index combinations
using checked equality reasoning. Reconstruct the eliminator application and
its relation to the original goal.
\end{enumerate}

Use Lean's own dependent-elimination machinery behind a wrapper rather than
reimplementing all details from constructor names. The wrapper should return a
mapping between old and new locals, the generated branch goals, and any
transport obligations. This mapping is needed by contracts and example
propagation, not only by pretty printing.

Case analysis is expensive when performed indiscriminately. Favor scrutinees
that occur in the target indices, appear in a discriminating contract, or
block reduction of a useful expression. Repeatedly splitting an irrelevant
Boolean is a valid but poor search move. A costed fallback can eventually try
other scrutinees without treating a heuristic relevance score as a theorem.

\subsection{Definitional equality versus proved equality}

Suppose an expression has type $V(A,n+m)$ while the goal is $V(A,m+n)$. A proof
that addition is commutative does not make these expressions definitionally
equal. The engine must either choose a differently oriented construction or
build a transport using a proof of $n+m=m+n$.

An equality-repair rule should first infer the exact source and target types.
If conversion fails, it searches for a small common family $F$ and indices
$i,j$ with source $F(i)$ and target $F(j)$. It then asks for $q:i=j$ and produces
the appropriate equality recursor applied to the source term. The motive is
checked in Lean; a raw cast operation is never introduced as a shortcut.
General heterogeneous equality should be an explicit extension, not an excuse
to replace unrelated types.

Arithmetic solvers are especially useful here. Their output must be a proof of
the needed index equation in the actual local context. A normalization that
uses truncated natural subtraction, integer division, or machine arithmetic
must match the source operation exactly. ``Both sides look like integers'' is
not a sufficient translation contract.

\subsection{Indices guide construction without specifying all behavior}

The prototype constructs \code{Fin n} from local $k,n$ and $k<n$, and constructs
a user-defined vector of length $n+1$ from a head and a length-$n$ tail. These
are direct uses of dependent constructor telescopes, not recognition of a
special built-in vector type.

Nevertheless, a function
\[
 (xs:\operatorname{List}(A))\to\operatorname{Fin}(|xs|)\to A
\]
does not by itself state that the selected element has the given index. On a
nonempty list, repeatedly returning the head can satisfy this type. Indexed
types remove invalid states and create useful obligations; behavioral
contracts still distinguish intended algorithms from other inhabitants.
The same distinction explains why a vector head result of type $A$ alone does
not document every desired law of an indexing API.

\subsection{Proof elimination and quotient-aware extensions}

Elimination must respect the declaration's actual recursor and result-sort
restrictions. From a proof of existence in \code{Prop}, the engine must not
invent an executable witness unless a permitted construction principle
supports it. Equality and other permitted eliminations should be handled by
their real rules rather than by a blanket ban on all proof-to-data dependence.

Quotient operations should initially be synthesized through registered
constructors and lifting APIs. A lift's well-definedness condition becomes a
proof obligation about the exact proposed function. Searching for arbitrary
quotient representatives is a different, often noncomputable request. This
API-oriented approach extends practical coverage without pretending that every
new semantic feature is just another algebraic data constructor.

\section{Recursion should be generated from types and contracts}\label{sec:recursion}

A broad practical synthesizer cannot depend solely on finding the right
nonrecursive composition. Lists, trees, syntax, finite maps, arrays with bounds,
and well-founded procedures recur throughout Lean developments. Recursion
therefore needs a dedicated planner, but it should generate \emph{typed
schemas from declarations}, not a fixed menu of named algorithms such as map,
filter, and reverse.

\subsection{Three routes, in a deliberate order}

First try a permitted library operation with a matching specification. Second
try a nonrecursive eliminator or fold already available in the local context.
Third generate a recursive schema from the input type. These routes can run in
parallel lanes, but this order is a useful default: reusing a correct library
implementation is often the right engineering answer, while provider-free mode
remains essential for evaluating synthesis itself.

A Church-encoded value may supply its own eliminator. Its carrier need not be a
simple data type: an index lookup can require a function-valued carrier such as
$\mathbb Z\to A$. The planner should derive carrier candidates from the demanded
result, remaining inputs, and contract rather than enumerate only base carriers.
This directly targets a class of composition difficulties visible in the
current repositories, without claiming that a native host alone solves their
original queries~\cite{indexing-diagnostic}.

\subsection{Structural schemas}

For an inductive input $I$, read the recursor's parameters, indices, major
premise, motive, minor premises, and recursive hypotheses. Choose a generalized
result family $M$. The schema's branch holes have the types of the instantiated
minor premises; recursive results are available exactly where the recursor
supplies them.

For lists and a nondependent result $B$, this yields the familiar shape
\[
 \operatorname{fold}(b,s,[])=b,\qquad
 \operatorname{fold}(b,s,x::xs)=s(x,\operatorname{fold}(b,s,xs)),
\]
but the engine should discover this shape from the recursor rather than encode
``a list has one recursive tail'' as a universal rule. A binary tree supplies
two recursive results. An indexed family changes the branch indices. More
complex nested or mutual recursors require their actual metadata and may be
staged as a later capability.

When generating surface recursive definitions instead of explicit recursor
terms, expose recursive calls through capabilities of the form
\[
 \mathsf{recCall}(x',\rho),
\]
where $\rho$ witnesses the selected decrease discipline. The branch grammar
must not contain an unrestricted local assumption $f:T$ representing the very
function being defined. Otherwise a proof search can ``solve'' the goal by
calling itself on the original arguments.

Lean ultimately justifies accepted structural recursion through recursors, and
well-founded recursion through a different termination construction. The
resulting computational equalities have different reduction behavior~\cite{lean-recursion}.
Leant2 should record which route it chose and generate proof obligations using
the appropriate equation lemmas, rather than assume every recursive equation
is reducible by reflexivity.

\subsection{Proof-carrying motives: compute the value and its invariant}

For a pointwise relation $Q(xs,ys)$, choose the motive
\begin{equation}\label{eq:packed-motive}
 M(xs)=\{ys:\operatorname{List}(B)\mid Q(xs,ys)\}.
\end{equation}
A recursive branch then receives both a smaller result $r$ and its proof
$Q(xs,r)$. The branch synthesizer constructs the next result and discharges its
local contract. Finally, project the data from the recursor result. The global
correctness theorem follows from the proof field.

For length preservation, the cons branch contains
\[
 r:\operatorname{List}(A),\qquad h:|r|=|xs|,
\]
and must produce $ys$ with $|ys|=|x::xs|$. Both $x::r$ and $r\mathbin{++}[x]$
meet that obligation. A stronger order-sensitive contract decides between
identity and reverse. Thus the motive helps prune but does not invent intent.

This is more effective than first completing an arbitrary recursive function
and only afterward attempting an unrelated induction proof. It makes the
induction hypothesis available at the point where the branch is chosen, and
contract failure can backtrack into that choice. It also avoids an unsound
circular proof: the hypothesis is supplied by a recursor on smaller data, not
assumed for the full current call.

For dependent output types, use the corresponding dependent result package or
an equivalent pair of computational/proof motives supported by the recursor.
For a function-valued output, the motive may include remaining arguments:
\[
 M(xs)=\prod_{a:S}\{y:B(xs,a)\mid Q(xs,a,y)\}.
\]
This is where generalizing an accumulator or an index becomes a substantive
algorithmic choice.

\subsection{Generalization and accumulator invention}

A fixed-result recursion schema may be too weak or inefficient. Generate a
bounded family of generalized motives by moving selected nondecreasing
arguments under the motive, introducing a candidate accumulator, or pairing
several mutually useful summaries. Prefer accumulator types already present
in the request or library specifications before inventing unrelated types.

For reversal, a useful generalized invariant is
\begin{equation}\label{eq:reverse-invariant}
 \operatorname{go}(xs,a)=\operatorname{reverse}(xs)\mathbin{++}a.
\end{equation}
The base branch returns $a$. The step recurses on $xs$ with accumulator $x::a$.
The induction hypothesis and list equations reduce the required branch law to
associativity and the cons/append identities. The public result is
$\operatorname{go}(xs,[])$.

The planner can propose this invariant from an expression grammar built from
the public relation, accumulator variable, available operations, and known
unit elements. Every proposed invariant has initiation, preservation, and
postcondition obligations. Failed induction is a failed invariant candidate,
not evidence that the desired function is impossible.

Accumulator invention should be costed separately from branch synthesis.
Otherwise the engine may endlessly add useless state fields. Limit the number
of new fields, their type-expression size, and the grammar of invariant
formulas per cost tier. Learned successful schemas can be generalized and
registered only with replayable proofs or as explicitly heuristic suggestions.

\subsection{Well-founded recursion}

Some natural programs recurse on a transformed input rather than a syntactic
subterm. Use a schema based on a well-founded relation $\prec$:
\[
 F(x)=\mathsf{step}\bigl(x,\lambda y\,(h:y\prec x).F(y)\bigr).
\]
The second argument of \code{step} only permits recursive calls when the branch
can prove the decrease. A proof-carrying variant returns both $F(x)$ and its
pointwise invariant, making the smaller-call hypotheses available to branch
proofs just as in~\eqref{eq:packed-motive}.

The initial measure grammar should include natural-valued input parameters,
registered size functions, differences with known bounds, and lexicographic
tuples. Generated measure candidates produce well-foundedness and per-call
decrease obligations. Arithmetic automation can prove many of these, but an
SMT guess that a quantity decreases is not enough. Measure synthesis is a search
component; termination is an acceptance obligation.

For example, an array loop incrementing $i$ can use $\operatorname{size}(a)-i$
under a bound invariant, whereas searching for a syntactically smaller $i$ is
misdirected. The array's mutation semantics and loop postcondition remain
separate obligations. A purely functional model should be available before
attempting an optimized mutable implementation.

\subsection{Mutual recursion, nested data, and helper lemmas}

Mutual recursion should be handled as a jointly synthesized strongly connected
component with a common termination argument. Each component's contract may
refer to smaller calls in the other components only through the generated
capabilities. A collection of separately proved-looking cyclic declarations is
not a correctness argument.

For nested datatypes, the recursor can supply higher-order induction hypotheses
or require auxiliary traversals. Begin with patterns for which the recursor
adapter has explicit tests; report an unsupported schema rather than lowering
the type to a simpler tree. This staged coverage still leaves library-based
construction available for those types.

Helper-lemma synthesis is valuable when a contract proof repeatedly gets stuck
at the same algebraic or relational shape. Propose a generalized lemma whose
parameters are abstracted from the failing goal, then prove it in a context
that does not assume the desired global result. Its dependency graph must be
acyclic, or justified by the same legal recursive construction. Unproved
lemmas cannot become providers merely because they make branch synthesis easy.

\subsection{Relational recursion and effects}

Injectivity or a two-run relation may require induction on two inputs or a
simulation relation between execution states. Extend the schema registry with
relational recursors and product-program proof rules rather than force such
contracts into a single-run postcondition. A structural program can be found
by the term lane while the proof lane chooses the right induction variables
and generalized relation.

For an effectful computation, choose a registered semantics such as a state
transformer, an exception result, or a weakest-precondition interpretation.
A \code{bind} rule can propagate a postcondition backward through the
continuation. An inhabitant of \code{IO A} alone says little about its side
effects. The default evaluator must not execute arbitrary effectful candidates
while testing their behavior.

Partial-fixpoint and coinductive facilities can be supported through separate
schema families with their own correctness meanings. They should not be
silently substituted when a total program's termination proof fails. This is
an extension of the request profiles, not a relaxation of
Equation~\eqref{eq:acceptance}.

\section{Library-directed synthesis and proof summaries}\label{sec:library}

\subsection{A provider is an exact term plus useful theorems}

Store each provider's exact declaration or local term, universe telescope,
binder structure, result-head indexes, instance dependencies, and permitted
use policy. A specification entry references a theorem whose type relates
that exact operation to its inputs and output. Registration checks the
reference and its scope. A human-written annotation is not accepted as an
axiom that the provider has some convenient semantics.

For instance, a list append provider can carry length, element, and order
properties. If the current contract only mentions length, the length theorem
can discharge a substantial part of the branch without unfolding append's
implementation. If the contract mentions order, a multiset summary is not a
replacement. Multiple specifications for the same provider can be selected
according to the demanded observations.

Separate \emph{term retrieval} from \emph{proof retrieval}. A theorem whose
statement mentions the target function may help prove an emitted composition;
it does not automatically grant permission to call the reference function in
the implementation. Likewise, a namespace blacklist is not a robust way to
hide an oracle if the same operation is reachable through an alias. Provider
exclusion should follow declaration dependencies and the explicit inventory.

\subsection{Backward contract propagation}

Suppose the sketch is $g(a)$, the desired postcondition is $Q(g(a))$, and a
registered theorem establishes
\[
 \forall a,\ P(a)\to Q(g(a)).
\]
Then $P(a)$ is a sufficient obligation for the argument synthesizer. This can
be much smaller than completing arbitrary arguments and evaluating them later.
However, sufficiency alone does not justify deleting all arguments that fail
$P$. They might satisfy $Q(g(a))$ for another reason. Keep this as one
construction lane unless an equivalence or a necessary-condition theorem
supports stronger pruning.

For a conditional, split the postcondition into branch conditions under the
exact guard assumptions. For a pattern match, use the branch substitutions
from Section~\ref{sec:dependent}. For a constructor, propagate observations to
fields using constructor laws. For an opaque higher-order provider, stop with
a symbolic residual unless a suitable theorem is registered.

This formulation is inspired by refinement-type synthesis and bidirectional
example propagation~\cite{synquid,smyth}. Leant2's distinguishing constraint is
that these rules operate on Lean's actual dependent expressions, and every
claimed implication must be reconstructed in its original logic.

\subsection{Law-carrying structures}

Many useful Lean types ask for operations and proofs of their laws together.
A monoid-like structure is not just three independent goals for multiplication,
a unit, and associativity. The proof fields constrain the same operation
metavariables; alternative units and multiplication terms must remain coupled.
Use the structure telescope to generate holes, then let cheap law checks reject
bad operation choices early.

Operations can also be transported along an equivalence. Register a generic
transport construction with its law proofs, and search for an appropriate
equivalence as a subproblem. This reuses a mathematically meaningful algorithm
instead of re-proving all laws from scratch. The equivalence itself is an exact
value in the term graph, including inverse laws, not a Boolean claim from a
plugin.

Specification precision remains important. For a stable sort, sortedness and a
permutation statement may not express preservation of the order of equal-key
records. For a filter, a set-level membership equivalence can lose multiplicity
information. The registry should support sequence- and position-sensitive
relations when requested, not pretend every behavioral contract reduces to
sets or lengths.
